\uuid{GKE8}
\titre{Limite en l'infini}
\auteur{bodin ; hanani ; mzari}
\organisation{exo7}
\licence{CC BY-NC-SA 4.0}
\datecreate{2020-01-20}
\niveau{L1}
\module{Analyse}
\chapitre{Continuité, limite et étude de fonctions réelles}
\sousChapitre{Limite de fonctions}
\difficulte{2}
\theme{Limites des fonctions réelles, limite en l'infini}
\exo7qcmid{lille-1.limites.016}

\contenu{
\begin{question}
\qclassification{123.03}
\qtags{motcle=limite en l'infini}
Soit $f(x)= \sin x$. Quelles sont les assertions vraies ?
\begin{answers}

     \bad{$\lim_{x\to +\infty} f(x)=1$}

    \good{$f$ n'admet pas de limite en $+\infty$.}

    \bad{$\lim_{x\to -\infty} f(x)=-1$}

    \good{$f$ n'admet pas de limite en $-\infty$.}

\end{answers}
\begin{explanations}
Toute fonction périodique non constante n'admet pas de limite en l'infini.
\end{explanations}
\end{question}
}
